\chapter{The final statements and the assembly}\label{ch:fin}

The archive's file \arch{Main.lean} states five theorems about Mathlib's
Riemann zeta function and Mathlib's Dirichlet $L$-functions. Each is a one-line
term built from a single nonvanishing statement: no Dirichlet $L$-function
vanishes in the half plane $\Re s>1/2$, with the pole of $\zeta$ excluded. This
chapter transcribes the five statements, the elementary deductions that turn the
nonvanishing statement into the other four, and the short chain of implications
by which the archive derives the nonvanishing statement. The chain ends at the
windowed analytic criterion, the supply of host fields, the Chebyshev inequality
on windows, the Tate--Shafarevich budget \texttt{VWin}, and the existence of
geometric families, which later chapters treat.

\section{The final statements}

\begin{definition}[Mathlib's Riemann hypothesis]\label{fin:rh-def}
\lean{RiemannHypothesis, RiemannZetaZerosOnLine, Claim.Zeta}
Let $\zeta$ denote Mathlib's \texttt{riemannZeta}, a function
$\CC\to\CC$ that equals the analytic continuation of the Riemann zeta function
on $\CC\setminus\{1\}$ and has some value at $s=1$. The \emph{Riemann
hypothesis} is the statement: for every $s\in\CC$ with $\zeta(s)=0$, with
$s\ne -2(n+1)$ for all $n\in\Z_{\ge0}$, and with $s\ne1$, one has $\Re s=1/2$.
The archive's \texttt{RiemannZetaZerosOnLine} (in \arch{Statements.lean}) and
\texttt{Claim.Zeta} have the same body as Mathlib's \texttt{RiemannHypothesis}.
\end{definition}

\begin{lemma}[the three forms agree]\label{fin:zeta-iff}
\lean{Claim.zeta_iff}
\uses{fin:rh-def}
$\texttt{Claim.Zeta}\iff\texttt{RiemannHypothesis}$.
\end{lemma}
\begin{proof}
\uses{fin:rh-def}
The two propositions are definitionally equal; the archive's proof is
reflexivity. The same holds for \texttt{RiemannZetaZerosOnLine}, which
\arch{Main.lean} proves with the term that proves \texttt{RiemannHypothesis}.
\end{proof}

\begin{definition}[nonvanishing right of $1/2$]\label{fin:every-character}
\lean{DirichletLNonvanishingRightOfHalf, UpperTheory.EveryCharacter, Claim.HalfPlane}
For $N\ge1$ and a Dirichlet character $\chi$ modulo $N$ (with values in $\CC$),
let $L(s,\chi)$ denote Mathlib's \texttt{DirichletCharacter.LFunction}, the
meromorphic continuation of $\sum_{n\ge1}\chi(n)n^{-s}$, which is entire unless
$\chi$ is the trivial character, and has a junk value at $s=1$ in that case. The
statement $\mathrm{EC}$ is: for every $N\ge1$, every Dirichlet character $\chi$
modulo $N$ and every $s\in\CC$ with $\Re s>1/2$ and ($\chi\ne1$ or $s\ne1$),
$L(s,\chi)\ne0$. The three Lean names have literally the same body.
\end{definition}

\begin{definition}[zeros in the strip]\label{fin:strip-def}
\lean{DirichletLZerosOnLine, Claim.Strip}
The statement: for every $N\ge1$, every Dirichlet character $\chi$ modulo $N$
and every $s$ with $0<\Re s<1$ and $L(s,\chi)=0$, one has $\Re s=1/2$.
\end{definition}

\begin{definition}[zeros in the strip, primitive characters]\label{fin:prim-def}
\lean{PrimitiveDirichletLZerosOnLine}
The statement: for every $N\ge1$, every \emph{primitive} Dirichlet character
$\chi$ modulo $N$ and every $s$ with $0<\Re s<1$ and $L(s,\chi)=0$, one has
$\Re s=1/2$.
\end{definition}

\begin{remark}[fidelity]\label{fin:rem-fidelity}
The file \arch{Statements.lean} imports only Mathlib and uses only Mathlib's
\texttt{riemannZeta}, \texttt{DirichletCharacter}, \texttt{LFunction} and
\texttt{IsPrimitive}. Mathlib's \texttt{RiemannHypothesis} was read in a Mathlib
checkout of a different revision from the pinned one
(\arch{Mathlib/NumberTheory/LSeries/RiemannZeta.lean}); its body there matches
the one transcribed above. The statement \texttt{RiemannZetaZerosOnLine} excludes
exactly the trivial zeros and the point $1$; the nonvanishing statement
excludes only the pair (trivial character, $s=1$), for every modulus $N$. Since
$L(s,\chi)$ for the trivial character modulo $N$ is $\zeta(s)\prod_{p\mid N}(1-p^{-s})$,
the nonvanishing statement at $\chi=1$ is the Riemann hypothesis in the form
``$\zeta(s)\ne0$ for $\Re s>1/2$, $s\ne1$''.
\end{remark}

\begin{theorem}[the five theorems of \arch{Main.lean}]\label{fin:main}
\lean{riemannHypothesis, riemannZeta_eq_zero_re, LFunction_ne_zero_of_half_lt_re, LFunction_eq_zero_re, LFunction_eq_zero_re_of_isPrimitive}
\uses{fin:rh-def, fin:every-character, fin:strip-def, fin:prim-def}
The following hold.
\begin{enumerate}
\item The Riemann hypothesis (\texttt{riemannHypothesis}), and the same
statement in the form \texttt{RiemannZetaZerosOnLine}
(\texttt{riemannZeta\_eq\_zero\_re}).
\item The nonvanishing statement $\mathrm{EC}$
(\texttt{LFunction\_ne\_zero\_of\_half\_lt\_re}).
\item The zeros of every $L(s,\chi)$ in $0<\Re s<1$ lie on $\Re s=1/2$
(\texttt{LFunction\_eq\_zero\_re}).
\item The same for primitive $\chi$ (\texttt{LFunction\_eq\_zero\_re\_of\_isPrimitive}).
\end{enumerate}
\end{theorem}
\begin{proof}
\uses{fin:kernels}
Each item is one of the kernel terms: (1) is \texttt{rhOfAll} (used twice, at
two definitionally equal types), (2) is \texttt{allChi}, (3) is the second
component of \texttt{claimAll}, and (4) is \texttt{critAll}.
\end{proof}

\begin{theorem}[the kernel terms]\label{fin:kernels}
\lean{allChi, rhOfAll, critAll, claimAll}
\uses{fin:every-character, fin:rh-def, fin:strip-def, fin:prim-def}
The statement $\mathrm{EC}$ holds (\texttt{allChi}); the Riemann hypothesis holds
(\texttt{rhOfAll}); the primitive strip statement holds (\texttt{critAll}); and
the conjunction of $\mathrm{EC}$, the strip statement and \texttt{Claim.Zeta}
holds (\texttt{claimAll}).
\end{theorem}
\begin{proof}
\uses{fin:every-character-thm, fin:rh-of-every, fin:crit-line, fin:strip, fin:zeta}
In \arch{Supplements/Kernels/MainTheorems.lean}: \texttt{allChi} is the
archive's theorem $\mathrm{EC}$ (Theorem~\ref{fin:every-character-thm});
\texttt{rhOfAll} applies Lemma~\ref{fin:rh-of-every} to it; \texttt{critAll}
applies Lemma~\ref{fin:crit-line} to it; \texttt{claimAll} is the triple
$(\texttt{allChi},\ \text{Lemma~\ref{fin:strip}}(\texttt{allChi}),\ \text{Lemma~\ref{fin:zeta}}(\texttt{allChi}))$.
All of the content therefore lies in the single statement $\mathrm{EC}$.
\end{proof}

\section{From nonvanishing to the other statements}

The deductions in this section are elementary consequences of Mathlib's
functional equations and Euler products; their Lean proofs were read in full.

\begin{lemma}[Riemann hypothesis from nonvanishing on the right]\label{fin:rh-of-right}
\lean{Formalism.VW.rh_of_right}
\uses{fin:rh-def}
Assume $\zeta(s)\ne0$ for every $s$ with $\Re s>1/2$ and $s\ne1$. Then the
Riemann hypothesis holds.
\end{lemma}
\begin{proof}
\uses{fin:rh-def}
Let $\zeta(s)=0$ with $s$ not a trivial zero and $s\ne1$. If $\Re s>1/2$ the
hypothesis is contradicted. Suppose $\Re s<1/2$. Then $s\ne0$ because
$\zeta(0)=-1/2$, so $1-s\ne1$, and $\Re(1-s)>1/2$, so $1-s$ is not a
nonpositive integer. Mathlib's functional equation \texttt{riemannZeta\_one\_sub},
applied at $1-s$, gives
\[
\zeta(s)=2\,(2\pi)^{-(1-s)}\,\Gamma(1-s)\,\cos\!\big(\pi(1-s)/2\big)\,\zeta(1-s).
\]
The power of $2\pi$ is nonzero, $\Gamma(1-s)\ne0$ since $\Re(1-s)>0$, and
$\zeta(1-s)\ne0$ by hypothesis. If $\cos(\pi(1-s)/2)=0$ then $s=-2k$ for an
integer $k$: $k\le-1$ contradicts $\Re s<1/2$, $k=0$ gives $s=0$, and $k\ge1$
makes $s$ a trivial zero. Hence the right side is nonzero, a contradiction. So
$\Re s=1/2$.
\end{proof}

\begin{lemma}[the Riemann hypothesis from EC]\label{fin:rh-of-every}
\lean{Formalism.VW.rh_of_every}
\uses{fin:every-character, fin:rh-def}
Assume $\mathrm{EC}$. Then the Riemann hypothesis holds.
\end{lemma}
\begin{proof}
\uses{fin:rh-of-right}
Apply $\mathrm{EC}$ with $N=1$ and $\chi$ the unique character modulo $1$, using
$s\ne1$. Mathlib's \texttt{LFunction\_modOne\_eq} identifies $L(s,\chi)$ with
$\zeta(s)$. Conclude with Lemma~\ref{fin:rh-of-right}.
\end{proof}

\begin{lemma}[gamma factors on the right half plane]\label{fin:gamma}
\lean{Formalism.VW.gammaFactor_ne_zero, Claim.gamma_ne_zero}
For every Dirichlet character $\chi$ and every $s$ with $\Re s>0$, the gamma
factor $\Gamma_{\R}(s+a_{\chi})$ of $\chi$ (Mathlib's \texttt{gammaFactor}, with
$a_{\chi}=0$ for even and $a_{\chi}=1$ for odd $\chi$) is nonzero. The second name
requires $N\ge1$.
\end{lemma}
\begin{proof}
Split by parity of $\chi$ and use $\Gamma_{\R}(t)\ne0$ for $\Re t>0$.
\end{proof}

\begin{lemma}[the critical line for primitive characters]\label{fin:crit-line}
\lean{Formalism.VW.critLine}
\uses{fin:every-character, fin:prim-def}
Assume $\mathrm{EC}$. Let $N\ge1$, let $\chi$ be a primitive character modulo
$N$, and let $0<\Re s<1$ with $L(s,\chi)=0$. Then $\Re s=1/2$.
\end{lemma}
\begin{proof}
\uses{fin:gamma}
Here $s\ne0,1$. If $\Re s>1/2$, $\mathrm{EC}$ (with $s\ne1$) is contradicted.
Suppose $\Re s<1/2$. Since $L(s,\chi)=\Lambda(s,\chi)/\Gamma_{\R}(s+a_{\chi})$ with
$\Lambda$ Mathlib's \texttt{completedLFunction}, and the gamma factor is
nonzero (Lemma~\ref{fin:gamma}), $\Lambda(s,\chi)=0$. The character $\chi^{-1}$
is primitive (its conductor equals that of $\chi$). Mathlib's functional
equation for primitive characters gives
$\Lambda(1-s,\chi^{-1})=\varepsilon\,\Lambda(s,\chi)=0$ for a root number
$\varepsilon$, so $L(1-s,\chi^{-1})=0$. As $\Re(1-s)>1/2$ and $1-s\ne1$, this
contradicts $\mathrm{EC}$ for $\chi^{-1}$ modulo $N$.
\end{proof}

\begin{lemma}[Euler factors]\label{fin:euler-factor}
\lean{Claim.factor_ne_zero}
Let $\psi$ be a Dirichlet character modulo some $M$, let $\Re s>0$ and let $p$
be prime. Then $1-\psi(p)p^{-s}\ne0$.
\end{lemma}
\begin{proof}
$|\psi(p)|\le1$ and $|p^{-s}|=p^{-\Re s}<1$.
\end{proof}

\begin{lemma}[passing to the primitive character]\label{fin:prim-zero}
\lean{Claim.prim_zero}
\uses{fin:euler-factor}
Let $N\ge1$, let $\chi$ be a character modulo $N$ and $\chi^{\ast}$ the primitive
character inducing it (Mathlib's \texttt{primitiveCharacter}, of level the
conductor of $\chi$). If $0<\Re s<1$ and $L(s,\chi)=0$, then $L(s,\chi^{\ast})=0$.
\end{lemma}
\begin{proof}
\uses{fin:euler-factor}
Since $s\ne1$, Mathlib's \texttt{LFunction\_changeLevel} gives
$L(s,\chi)=L(s,\chi^{\ast})\prod_{p\mid N}(1-\chi^{\ast}(p)p^{-s})$. The product is
nonzero by Lemma~\ref{fin:euler-factor}.
\end{proof}

\begin{lemma}[reflection]\label{fin:prim-reflect}
\lean{Claim.prim_reflect}
\uses{fin:gamma}
Let $\psi$ be primitive modulo $M\ge1$, let $0<\Re s<1$ and $L(s,\psi)=0$. Then
$L(1-s,\psi^{-1})=0$.
\end{lemma}
\begin{proof}
\uses{fin:gamma}
As in the proof of Lemma~\ref{fin:crit-line}: the completed $L$-function
vanishes at $s$, the functional equation for the primitive character
$\psi^{-1}$ moves the zero to $1-s$, and $1-s\ne0$ lets one divide by the gamma
factor again.
\end{proof}

\begin{lemma}[the strip statement]\label{fin:strip}
\lean{Claim.strip}
\uses{fin:every-character, fin:strip-def}
Assume $\mathrm{EC}$. Then the strip statement holds for all characters.
\end{lemma}
\begin{proof}
\uses{fin:prim-zero, fin:prim-reflect}
Let $L(s,\chi)=0$ with $0<\Re s<1$. If $\Re s>1/2$, $\mathrm{EC}$ is
contradicted. If $\Re s<1/2$, Lemma~\ref{fin:prim-zero} gives
$L(s,\chi^{\ast})=0$, Lemma~\ref{fin:prim-reflect} gives
$L(1-s,(\chi^{\ast})^{-1})=0$, and since $\Re(1-s)>1/2$ and $1-s\ne1$ this
contradicts $\mathrm{EC}$ at the modulus equal to the conductor of $\chi$.
\end{proof}

\begin{lemma}[the zeta statement]\label{fin:zeta}
\lean{Claim.zeta}
\uses{fin:every-character, fin:rh-def}
Assume $\mathrm{EC}$. Then \texttt{Claim.Zeta}, i.e.\ the Riemann hypothesis,
holds.
\end{lemma}
\begin{proof}
\uses{fin:strip}
Let $\zeta(s)=0$, $s$ not a trivial zero, $s\ne1$; then $s\ne0$. For
$0<\Re s<1$ apply Lemma~\ref{fin:strip} to the character modulo $1$. For
$\Re s\ge1$ use Mathlib's nonvanishing $\zeta(s)\ne0$ on $\Re s\ge1$. For
$\Re s\le0$: $\Gamma_{\R}(s)\ne0$ because $s$ is neither $0$ nor a trivial zero,
so the completed zeta function vanishes at $s$, hence at $1-s$ by its functional
equation, so $\zeta(1-s)=0$ with $\Re(1-s)\ge1$, a contradiction. This
component of \texttt{claimAll} is not used by \arch{Main.lean}.
\end{proof}

\begin{lemma}[a second derivation of the Riemann hypothesis]\label{fin:rh-of-everychar}
\lean{UpperTheory.zeta_ne_zero_of_everyCharacter, UpperTheory.riemannHypothesis_of_everyCharacter, UpperTheory.FinalTheorems.riemannHypothesis}
\uses{fin:every-character, fin:rh-def}
Assume $\mathrm{EC}$. Then $\zeta(s)\ne0$ for $\Re s>1/2$, $s\ne1$, and the
Riemann hypothesis holds. In particular the Riemann hypothesis holds
unconditionally in the archive (\texttt{UpperTheory.FinalTheorems.riemannHypothesis}).
\end{lemma}
\begin{proof}
\uses{fin:every-character-thm}
The first claim is $\mathrm{EC}$ modulo $1$. For the second, a zero with
$\Re s<1/2$ is reflected through the functional equation of the completed zeta
function $\xi(s)=\Gamma_{\R}(s)\zeta(s)$; the case $\Gamma_{\R}(s)=0$ forces $s=0$
or a trivial zero. The third name applies this to
Theorem~\ref{fin:every-character-thm}. None of these three declarations is used
by \arch{Main.lean}, which goes through Lemma~\ref{fin:rh-of-every}.
\end{proof}

\section{Primitive characters and the imprimitive step}





\begin{remark}[the trivial character at $s=1$]\label{fin:rem-junk}
$\mathrm{ZeroFreeAll}$ has no exception at $s=1$. For the trivial primitive
character it asserts $L(1,1)=\zeta(1)\ne0$, where $\zeta(1)$ is Mathlib's junk
value $(\gamma-\log4\pi)/2$, which is indeed nonzero. So $\mathrm{ZeroFreeAll}$
is slightly stronger than the primitive case of $\mathrm{EC}$, and whatever
proves it must, somewhere, bound the value Mathlib assigns at the pole. This was
not traced.
\end{remark}

\begin{lemma}[the imprimitive step]\label{fin:imprimitive}
\lean{UpperTheory.Top.imprimitive}
\uses{an:zerofree, fin:every-character}
Assume $\mathrm{ZeroFreeAll}$. Then $\mathrm{EC}$ holds.
\end{lemma}
\begin{proof}
\uses{an:pchar, fin:euler-factor}
Let $\chi$ be a character modulo $N\ge1$, $\Re s>1/2$, and ($\chi\ne1$ or
$s\ne1$). $\mathrm{ZeroFreeAll}$ at $\mathrm{of}(\chi)$ gives
$L(s,\chi^{\ast})\ne0$. If $\chi\ne1$ then $\chi^{\ast}\ne1$, since the trivial
character changes level to the trivial character; so ($\chi^{\ast}\ne1$ or
$s\ne1$), which is the hypothesis of Mathlib's \texttt{LFunction\_changeLevel}:
$L(s,\chi)=L(s,\chi^{\ast})\prod_{p\mid N}(1-\chi^{\ast}(p)p^{-s})$. Each factor is
nonzero as in Lemma~\ref{fin:euler-factor} (the archive repeats that argument
inline).
\end{proof}

\section{The assembly chain}

\begin{definition}[the literature record]\label{fin:open-lit}
\lean{UpperTheory.LiteratureTheorems.OpenLiteratureStatements}
% uses-lean: UpperTheory.One.M.Row, UpperTheory.Zero.LLC.RecTExists
\uses{fam:one-m-row, fam:rect-exists} % resolved
The proposition $\mathrm{OLS}$ is the conjunction of three fields:
\begin{enumerate}
\item \texttt{caraianiLocalGlobalDual}: the proposition \texttt{UpperTheory.One.M.Row};
\item \texttt{localLanglandsExists}: the proposition \texttt{UpperTheory.Zero.LLC.RecTExists};
\item \texttt{galoisRepOfPolarizedPairLit}: the proposition \texttt{UpperTheory.One.M.Row} again.
\end{enumerate}
Here \texttt{One.M.Row} is \texttt{UpperTheory.GalStar.Row0b} at the recognition
notion \texttt{IsRecT}: for every CM field $E$, prime $\ell$, isomorphism
$\overline{\Q}_{\ell}\cong\CC$, $a\ge1$, cuspidal $\pi$ of $\mathrm{GL}_a$ over $E$
and Hecke character $\chi$ of the maximal totally real subfield with $\pi$
regular algebraic and $(\pi,\chi)$ polarized, there is a Galois representation
$r:G_E\to\GL_a(\overline{\Q}_{\ell})$ with weights and Hodge--Tate data satisfying
the archive's property list \texttt{PropsB}. \texttt{RecTExists} says that at
every finite place of every number field there is a map from irreducible
representations of $\GL_n$ to Weil--Deligne representations satisfying the
archive's laws \texttt{LawsT}. Both are treated in the chapter on Galois
representations.
\end{definition}

\begin{remark}[the record is weaker than its names]\label{fin:rem-olit-names}
The field names suggest Caraiani's local-global compatibility, the local
Langlands correspondence, and the existence of Galois representations for
polarized pairs. In fact fields (1) and (3) have the same type, and field (2)
states only the existence of a map with the laws \texttt{LawsT}, with no
uniqueness and no bijectivity (the archive states uniqueness separately as
\texttt{RecTUnique}, which this record does not contain). Nothing in the record
is left open: Theorem~\ref{fin:every-character-of-row} fills all three fields
with in-archive theorems. A second, unrelated declaration named
\texttt{RecTExists} exists in
\arch{Development/AbelianVarietyTheory/Definitions/NonGaloisDescent/RecTExists.lean};
the one used here is in namespace \texttt{UpperTheory.Zero.LLC}.
\end{remark}

\begin{lemma}[where the record's fields are used]\label{fin:lit-fields}
\lean{UpperTheory.LiteratureTheorems.local_langlands_recT_exists_of_literature, UpperTheory.LiteratureTheorems.local_langlands_gl_recT_exists_of_literature, UpperTheory.LiteratureTheorems.pair_u_of_literature, UpperTheory.LiteratureTheorems.carU_local_global_of_literature}
\uses{fin:open-lit}
% uses-lean: UpperTheory.Zero.LLC.RecTExists, UpperTheory.One.M.New.llOf_alt, UpperTheory.GalS.field_of_starFull_alt, UpperTheory.GalS.ofDual_star_alt
\uses{fam:ll-of-alt, fam:rect-exists} % resolved
Assume $\mathrm{OLS}$. Then \texttt{RecTExists} holds (in two ways, the first the
field itself and the second obtained from it through \texttt{llOf\_alt} together
with a dozen in-archive published statements of local representation theory),
and the archive's propositions \texttt{PairUAlt} and \texttt{CarUAlt} hold.
\end{lemma}
\begin{proof}
\uses{fin:open-lit}
The four declarations project the fields of $\mathrm{OLS}$: the first returns
field (2); the second feeds it to \texttt{llOf\_alt}; \texttt{pair\_u} feeds
field (2) and field (3) to \texttt{field\_of\_starFull\_alt}; \texttt{carU\_local\_global}
feeds field (1) to \texttt{ofDual\_star\_alt}. Field (1) is also projected once
more, into the record \texttt{GlobalAutomorphicGaloisHypotheses'} built by
\texttt{literature\_automorphy}. These are the only direct projections of
$\mathrm{OLS}$ in \arch{.../ThirdPlaceSeries/EveryCharacterOfLiteratureOfVWin.lean}.
\end{proof}

\begin{remark}[the of-literature file]\label{fin:rem-plumbing}
The file \arch{.../ThirdPlaceSeries/EveryCharacterOfLiteratureOfVWin.lean}
contains $195$ theorems, each taking $\mathrm{OLS}$ as argument. Most return an
in-archive theorem and ignore the argument; for instance
\texttt{deligne\_ribet\_of\_literature} is the archive's
\texttt{PositivityTheory.PfDR.deligneRibet\_holds}. Only the four declarations
of Lemma~\ref{fin:lit-fields} project a field directly; others reach the fields
by calling those. The file carries no mathematics of its own.
\end{remark}

\begin{lemma}[the literature records]\label{fin:lit-records}
\lean{UpperTheory.LiteratureTheorems.published_results_eigen, UpperTheory.LiteratureTheorems.published_results_first, UpperTheory.LiteratureTheorems.theta_integrals_agree_and_local_matching_of_literature}
\uses{fin:open-lit, fin:lit-fields}
% uses-lean: GlobalTheory.Fc.S0.UnitaryGroupThirdPlaceHypotheses, GlobalTheory.FunPr.ThetaLocalInputs, GlobalTheory.IndefiniteRowSideFullHecke.WeilTensorModels.ThetaIntegralsAgree, GlobalTheory.IndefiniteRowSideFullHecke.WeilTensorModels.LocalMatching
Assume $\mathrm{OLS}$. Then the hypothesis records
\texttt{UnitaryGroupThirdPlaceHypotheses} (about $60$ fields named after
published theorems of automorphic representation theory, among them Green's
regular characters and the record \texttt{nLink}) and \texttt{ThetaLocalInputs}
hold, and \texttt{ThetaIntegralsAgree} and \texttt{LocalMatching} hold.
\end{lemma}
\begin{proof}
\uses{fin:lit-fields}
% uses-lean: PositivityTheory.ThKR.theoremK_head2
Each field of each record is filled by an of-literature theorem of the same
file. The pair \texttt{ThetaIntegralsAgree}, \texttt{LocalMatching} comes from
\texttt{theoremK\_head2} applied to two further of-literature theorems
(\texttt{eis\_bl}, \texttt{siegel\_weil\_std}). The field
\texttt{localGlobalAboveL} uses \texttt{carU\_above\_l\_of\_literature}, which uses
Lemma~\ref{fin:lit-fields}.
\end{proof}

\begin{theorem}[EC from the green parametrisation, VWin and a hypothesis record]\label{fin:every-green}
\lean{UpperTheory.Top.zero_free_all_of_green_v_win, UpperTheory.Top.every_character_of_green_v_win}
\uses{an:zerofree, fin:every-character}
% uses-lean: GlobalTheory.GLq.Green, UpperTheory.Top.VWin, GlobalTheory.Fc.HypsAlt, GlobalTheory.Fc.DefiniteSideCondition, Formalism.MainNotions.Closed, Formalism.MainNotions.LinkN
\uses{fam:green, fam:hyps-alt, geob:closed, sha:vwin} % resolved
Assume \texttt{Green} (the archive's parametrisation statement, treated with the
families), $\mathrm{VWin}$ (Tate--Shafarevich budgets on windows with constant
$401$ for every primitive character), and that the hypothesis record
$\texttt{HypsAlt}(N,\texttt{DefiniteSideCondition})$ is inhabited, where
$N(\mathcal C)$ says that the link predicate \texttt{LinkN} is closed on the
class of families $\mathcal C$. Then $\mathrm{ZeroFreeAll}$ holds, and hence
$\mathrm{EC}$ holds.
\end{theorem}
\begin{proof}
\uses{fin:imprimitive}
% uses-lean: UpperTheory.Top.zeroFreeAll_of_topG, GlobalTheory.Fc.forall_goal_of_shape_of_hyps, UpperTheory.Top.shape_closed_link_n_of_green, UpperTheory.Top.supplyChi, UpperTheory.Top.TopG, UpperTheory.Top.GChi, GlobalTheory.EX.HostExW
\uses{an:zerofreeall-topg, fam:goal-of-shape, fam:shape-closed, geob:topg, host:gchi, host:host-exw, host:supply-chi} % resolved
Take an element $h$ of the hypothesis record. The windowed analytic criterion
\texttt{zeroFreeAll\_of\_topG} with $\lambda=3/2$ needs two inputs. The first is
the Chebyshev inequality at window scale \texttt{TopG}; it is obtained from
\texttt{shape\_closed\_link\_n\_of\_green} applied to \texttt{Green}, which
produces the shape statement $\texttt{TopShapeAlt}(\text{true},\text{true},N,\texttt{TopG})$,
and \texttt{forall\_goal\_of\_shape\_of\_hyps}, which discharges that shape with
$h$. The second is $\texttt{HostExW}(3/2,G_{\psi})$ for every primitive $\psi$,
obtained from \texttt{supplyChi} applied to the $\psi$-component of
$\mathrm{VWin}$. This gives $\mathrm{ZeroFreeAll}$; the second name then applies
Lemma~\ref{fin:imprimitive}.
\end{proof}

\begin{theorem}[EC from the theta-integral records]\label{fin:theta-local}
\lean{UpperTheory.Lit.every_character_of_theta_integrals_local_matching}
\uses{fin:every-character}
% uses-lean: GlobalTheory.Fc.S0.UnitaryGroupThirdPlaceHypotheses, GlobalTheory.FunPr.ThetaLocalInputs, GlobalTheory.IndefiniteRowSideFullHecke.WeilTensorModels.ThetaIntegralsAgree, GlobalTheory.IndefiniteRowSideFullHecke.WeilTensorModels.LocalMatching, GlobalTheory.CovTop.CovolumeBoundOfX, UpperTheory.Zero.LLC.RecTExists, UpperTheory.Top.VWin
\uses{fam:rect-exists, geob:covolbound, sha:vwin} % resolved
Assume the records \texttt{UnitaryGroupThirdPlaceHypotheses} ($P$) and
\texttt{ThetaLocalInputs} ($F$), the statements \texttt{ThetaIntegralsAgree},
\texttt{LocalMatching}, the covolume bound \texttt{CovolumeBoundOfX},
\texttt{RecTExists} and $\mathrm{VWin}$. Then $\mathrm{EC}$ holds.
\end{theorem}
\begin{proof}
\uses{fin:every-green}
% uses-lean: GlobalTheory.GreenBridge.green_of_regular, UpperTheory.Lit.nonempty_hyps_definite_side_of_closed_bound, UpperTheory.N.boundAtTwoThree_linked_of_recT, GlobalTheory.Mem.Top.Bl.parametersOfDependence
\uses{fam:nonempty-hyps} % resolved
Apply Theorem~\ref{fin:every-green}. \texttt{Green} comes from
\texttt{green\_of\_regular} applied to the field
\texttt{greenRegularCharacters} of $P$. The hypothesis record comes from
\texttt{nonempty\_hyps\_definite\_side\_of\_closed\_bound}, applied to all seven
inputs except $\mathrm{VWin}$ and to the family of statements, indexed by
$c\in\R$, that \texttt{boundAtTwoThree\_linked\_of\_recT} produces from the field
\texttt{nLink} of $P$ at the parameters \texttt{parametersOfDependence} of the
local Langlands input at $c$. That construction is the archive's existence of
geometric families and is treated in the chapter on families.
\end{proof}

\begin{theorem}[EC from the literature record and VWin]\label{fin:every-lit-vwin}
\lean{UpperTheory.LiteratureTheorems.every_character_of_literature_of_vWin}
\uses{fin:open-lit, fin:every-character}
% uses-lean: UpperTheory.Top.VWin
\uses{sha:vwin} % resolved
Assume $\mathrm{OLS}$ and $\mathrm{VWin}$. Then $\mathrm{EC}$ holds.
\end{theorem}
\begin{proof}
\uses{fin:theta-local, fin:lit-records, fin:lit-fields}
% uses-lean: SelfIntersectionVolumeEstimates.topCovolSlim2_use_alt_alt
\uses{geob:covolsupply} % resolved
Apply Theorem~\ref{fin:theta-local}. The records $P$, $F$ and the pair of
theta statements come from Lemma~\ref{fin:lit-records}. The covolume bound is
\texttt{topCovolSlim2\_use\_alt\_alt} applied to fifteen of-literature theorems
of arithmetic geometry (resolution and principalisation in characteristic zero,
Poincar\'e--Lelong, arithmetic Hilbert--Samuel and Moriwaki bigness,
toroidal compactifications, Grothendieck deformation and others).
\texttt{RecTExists} is the second declaration of Lemma~\ref{fin:lit-fields}.
\end{proof}

\begin{theorem}[EC from a row]\label{fin:every-character-of-row}
\lean{UpperTheory.FinalTheorems.everyCharacter_of_row}
\uses{fin:open-lit, fin:every-character}
% uses-lean: UpperTheory.GalStar.Row0b, UpperTheory.Zero.LLC.IsRecT
\uses{fam:laws-t, fam:props-b} % resolved
Assume \texttt{Row0b} at the recognition notion \texttt{IsRecT}, i.e.\ the
proposition \texttt{One.M.Row} of Definition~\ref{fin:open-lit}. Then
$\mathrm{EC}$ holds.
\end{theorem}
\begin{proof}
\uses{fin:every-lit-vwin}
% uses-lean: UpperTheory.Zero.LLC.recTExists, UpperTheory.One.vwin_of_row
\uses{fam:rect-exists-thm, sha:vwin-of-row} % resolved
Let $h$ be the hypothesis. The record $\mathrm{OLS}$ is
$\langle h,\ \texttt{recTExists},\ h\rangle$: the row fills both row fields, and
the in-archive theorem \texttt{recTExists} fills the local Langlands field.
$\mathrm{VWin}$ is \texttt{vwin\_of\_row}$(h)$. Apply
Theorem~\ref{fin:every-lit-vwin}.
\end{proof}

\begin{theorem}[EC]\label{fin:every-character-thm}
\lean{UpperTheory.FinalTheorems.everyCharacter}
\uses{fin:every-character}
The statement $\mathrm{EC}$ holds.
\end{theorem}
\begin{proof}
\uses{fin:every-character-of-row}
% uses-lean: UpperTheory.Gal.Int.galoisLocalGlobalStrong_recParameter
\uses{fam:gal-strong} % resolved
Apply Theorem~\ref{fin:every-character-of-row} to the in-archive theorem
\texttt{galoisLocalGlobalStrong\_recParameter}, which proves \texttt{Row0b} at
the notion \texttt{recT}, the archive's name for $(h,\mathrm{rc})\mapsto\texttt{IsRecT}\,h\,\mathrm{rc}$.
\end{proof}

\begin{remark}[the chain, and where its content is]\label{fin:rem-chain}
Reading downward from \arch{Main.lean}:
$\text{Theorem~\ref{fin:main}}\leftarrow\text{Theorem~\ref{fin:kernels}}\leftarrow
\text{Theorem~\ref{fin:every-character-thm}}\leftarrow
\text{Theorem~\ref{fin:every-character-of-row}}\leftarrow
\text{Theorem~\ref{fin:every-lit-vwin}}\leftarrow
\text{Theorem~\ref{fin:theta-local}}\leftarrow
\text{Theorem~\ref{fin:every-green}}$, plus Lemma~\ref{fin:imprimitive} and
the elementary deductions of the second section. Every step above is a one-line
term or a short elementary proof. The mathematical content sits in the inputs
named in \texttt{uses-lean} comments: the windowed analytic criterion and its
\texttt{FormB} input, the host fields of \texttt{supplyChi} and
\texttt{HostExW}, the Chebyshev inequality \texttt{TopG} with
\texttt{level\_gm\_le\_on\_family\_scales}, the budget $\mathrm{VWin}$ from
\texttt{vwin\_of\_row}, and the row \texttt{One.M.Row}, \texttt{RecTExists} and
the family construction behind \texttt{shape\_closed\_link\_n\_of\_green} and
\texttt{nonempty\_hyps\_definite\_side\_of\_closed\_bound}.
\end{remark}

\begin{theorem}[an unused copy of the top of the chain]\label{fin:plug}
\lean{UpperTheory.Top.plugWin, UpperTheory.Top.plugAll, UpperTheory.Lit.everyE₂}
\uses{an:zerofree, fin:every-character}
% uses-lean: GlobalTheory.GLq.Green, UpperTheory.Top.VWin, GlobalTheory.Fc.Hyps, GlobalTheory.Fc.DRS
\uses{fam:green, fam:top-shape-old, sha:vwin} % resolved
Assume \texttt{Green}, $\mathrm{VWin}$, and that the record
$\texttt{Hyps}(N,\texttt{DRS})$ is inhabited, with $N$ as in
Theorem~\ref{fin:every-green}. Then $\mathrm{ZeroFreeAll}$ holds
(\texttt{plugWin}) and $\mathrm{EC}$ holds (\texttt{plugAll}). Further,
\texttt{everyE₂} derives $\mathrm{EC}$ from records \texttt{HypothesesUE},
\texttt{Hypotheses₂}, the two theta statements, \texttt{CovolX},
\texttt{LocalLanglandsGL} and $\mathrm{VWin}$.
\end{theorem}
\begin{proof}
\uses{fin:imprimitive}
% uses-lean: UpperTheory.Top.zeroFreeAll_of_topG, GlobalTheory.Fc.shape_of, UpperTheory.Top.shapeWin, UpperTheory.Top.supplyChi, UpperTheory.Lit.hypsE₂
\uses{an:zerofreeall-topg, fam:shape-win, host:supply-chi} % resolved
\texttt{plugWin} is the proof of Theorem~\ref{fin:every-green} with
\texttt{shapeWin} and \texttt{shape\_of} in place of
\texttt{shape\_closed\_link\_n\_of\_green} and
\texttt{forall\_goal\_of\_shape\_of\_hyps}; \texttt{plugAll} composes it with
Lemma~\ref{fin:imprimitive}; \texttt{everyE₂} applies \texttt{plugAll} to
\texttt{green\_of\_regular} and \texttt{hypsE₂}.
\end{proof}

\begin{remark}[the copy is not on the path from \arch{Main.lean}]\label{fin:rem-deadcopy}
A textual search of the archive finds \texttt{plugWin} used only by
\texttt{plugAll}, \texttt{plugAll} used only by \texttt{everyE₂}, and no use of
\texttt{everyE₂}. The path from \arch{Main.lean} goes through
\arch{.../ThirdPlaceSeries/TopPlug.lean} (Theorem~\ref{fin:every-green}) and
\arch{.../ThirdPlaceSeries/SiegelEisensteinSecondTermHypotheses.lean}
(Theorem~\ref{fin:theta-local}). The definition $\mathrm{EC}$ itself lives in
the file of the copy, \arch{Development/AutomorphicArithmetic/CharacterBoxes/EveryCharacter.lean}.
\end{remark}
